
\begin{obeassessment}
\textbf{Kuis Bab 3}
\begin{enumerate}
    \item Mengapa pemahaman tentang bilangan prima dan relatif prima sangat penting
          dalam algoritma kunci publik seperti RSA?
    \item Jika sebuah cipherteks memiliki nilai entropi yang sangat rendah, apa
          implikasinya terhadap keamanan pesan tersebut?
    \item Hitung $-37 \bmod 12$ dan $5^{-1} \bmod 12$. Tunjukkan langkah pemeriksaannya.
    \item Hitung $\varphi(21)$ dengan dua cara: mendaftar bilangan yang relatif prima
          terhadap 21, dan memakai sifat multiplikatif karena $21 = 3 \cdot 7$.
    \item Jelaskan mengapa keberadaan invers modulo $a^{-1} \bmod n$ bergantung pada
          nilai $\text{PBB}(a,n)$, dan buktikan dengan Identitas B\'ezout.
    \item Sebuah cipherteks 400 karakter memiliki entropi 4,55 bit/karakter. Berapa
          persen entropi maksimum alfabet 26 huruf yang dicapainya, dan mengapa angka
          itu belum cukup untuk menyimpulkan bahwa ciphernya aman?
    \item Mengapa menentukan $x$ pada $g^x \equiv h \pmod p$ dianggap sulit, padahal
          menghitung $g^x \bmod p$ dari $x$ yang diketahui sangat mudah?
\end{enumerate}
\end{obeassessment}
\begin{competencychecklist}
\begin{tabularx}{\linewidth}{|X|c|}
\hline
Indikator Kompetensi & Tercapai \\
\hline
Saya dapat melakukan operasi aritmetika modulo dan kekongruenan & [ ] \\
\hline
Saya dapat menentukan apakah dua bilangan relatif prima & [ ] \\
\hline
Saya dapat menghitung PBB dengan Algoritma Euclidean dan koefisien B\'ezout-nya & [ ] \\
\hline
Saya dapat menghitung invers modulo & [ ] \\
\hline
Saya dapat menghitung fungsi totient Euler dan menerapkan Teorema Euler & [ ] \\
\hline
Saya dapat menjelaskan akar primitif dan persoalan logaritma diskrit & [ ] \\
\hline
Saya dapat menjelaskan konsep entropi dan menghitung nilainya & [ ] \\
\hline
Saya dapat menafsirkan entropi cipherteks dan kaitannya dengan analisis frekuensi & [ ] \\
\hline
\end{tabularx}
\end{competencychecklist}
